%% Capítulo de la propuesta: CAPÍTULO III con subsecciones 3.1 a 3.4
%% (formato de la cátedra). El número de capítulo lo fija main.tex.
\chapter{Sistema web para la gestión de paquetes turísticos y reservas}
\label{cap:propuesta}

\section*{Antecedentes del capítulo}

Vizcacha Tours, agencia turística de Potosí, atiende solicitudes de viajes
por diferentes canales: la oficina, las llamadas telefónicas, los mensajes y
las redes sociales. La necesidad de disponer de información uniforme sobre
paquetes, fechas, cupos y reservas motivó el desarrollo de una solución web
que centraliza el proceso y apoya al personal en la confirmación de las
solicitudes.

\section{Marco referencial}

\textbf{Sistema de información transaccional (TPS).} Es un sistema orientado
al registro y procesamiento de las operaciones del negocio. En Vizcacha Tours
se aplica a las reservas, los cupos y los pagos, donde la consistencia de las
transacciones resulta esencial.

\textbf{Gestión de reservas turísticas.} Comprende la solicitud, la
disponibilidad, la asignación de plazas y la confirmación de servicios. Las
plataformas citadas en la introducción \parencite{amadeusGDS,ezus}
constituyen referencias funcionales, sin que se contemple integrarlas al
sistema local.

\textbf{Kanban y STATIK.} Kanban permite gestionar el flujo de trabajo
mediante la visualización, los límites de trabajo en curso y la mejora
continua; STATIK es el procedimiento de seis pasos para diseñar un sistema
Kanban ajustado al contexto del servicio \parencite{kanbanGuia}.

\textbf{Modelado UML.} El Lenguaje Unificado de Modelado permite representar
la estructura y el comportamiento del sistema mediante casos de uso, clases,
objetos, actividades, secuencias, componentes y despliegue
\parencite{omgUml}.

\textbf{Estimación de software.} COCOMO básico relaciona el tamaño estimado
del producto con el esfuerzo y la duración esperados \parencite{boehm1981}.
La estimación ágil emplea el tamaño relativo de las historias y el ritmo
observado en el propio trabajo.

\textbf{Seguridad y experiencia de usuario.} El diseño considera permisos por
rol, protección de credenciales, validación de datos y facilidad de uso, en
concordancia con las recomendaciones de OWASP Foundation
\parencite{owasp2025} y de ISO 9241-210 \parencite{iso9241}.

\subsection{Herramientas para el desarrollo}

Se utilizó Kanban como método de gestión del flujo de trabajo, con el tablero
del repositorio para visualizar las historias, limitar el trabajo en curso e
identificar bloqueos. Las herramientas tecnológicas empleadas se resumen en
la tabla~\ref{tab:herramientas}.

\begin{tablaapa}{tab:herramientas}{Herramientas de desarrollo}
\small
\begin{tabular}{@{}p{4.6cm}p{10.4cm}@{}}
\toprule
Herramienta & Aplicación \\
\midrule
Kanban (tablero del repositorio) & Backlog, tarjetas, responsables, límites
de trabajo en curso y seguimiento. \\
Visual Studio Code & Edición y depuración del código. \\
HTML, CSS y JavaScript & Interfaz web adaptable a celulares y
computadoras. \\
Nuxt 4 / Nitro (Node.js 22) & Lógica de negocio, autenticación y API del
sistema. \\
PostgreSQL (Neon) & Persistencia relacional y transacciones para controlar
los cupos. \\
GitHub & Control de versiones y revisión de cambios. \\
PlantUML & Modelado del proceso y del diseño técnico. \\
Pruebas de API y navegador & Pruebas unitarias, de integración y de
aceptación. \\
\bottomrule
\end{tabular}
\nota{Elaboración propia.}
\end{tablaapa}

\section{Estimación del esfuerzo y los costos}

La estimación se elaboró con dos enfoques complementarios. El primero,
COCOMO básico \parencite{boehm1981}, parte de un tamaño hipotético de 3.000
líneas de código y sirve como hipótesis de planificación. El segundo es la
estimación ágil, que emplea los datos reales del proyecto: las 19 historias
del backlog, con 55 puntos de historia y 54 tareas, y el ritmo medido en el
tablero durante la primera ventana de trabajo, de 4,00 horas-persona por
punto de historia. Con la tarifa declarada de 5.000 bolivianos por
mes-persona, el alcance completo equivale a 1,38 mes-persona y 6.875 Bs, de
los cuales 2.500 Bs corresponden al producto mínimo viable de la primera
ventana (HU-08, HU-01 y HU-04). Las tareas de corrección identificadas
durante las pruebas se registraron dentro de las historias correspondientes,
de modo que el ritmo medido las absorbe sin requerir una estimación
separada; la proyección se recalcula al cierre de cada ventana. Los
supuestos, los cálculos y el resumen comparativo se presentan en el Anexo~A.

\section{Descripción del proceso de negocio}

\subsection{Diagrama de entrada y salida}

Las entradas provienen de los usuarios y del personal de la agencia. El
sistema procesa los datos de manera centralizada, verifica los permisos y la
disponibilidad, y genera los resultados que respaldan la atención al cliente
(figura~\ref{fig:ips-negocio}).

\begin{figuraapa}{fig:ips-negocio}{Diagrama de entrada, proceso y salida del proceso de reservas}
  \incluirDiagrama{procesos/negocio-entrada-proceso-salida}
  \nota{Fuente: elaboración propia.}
\end{figuraapa}

El cliente consulta paquetes y disponibilidad por internet o en la oficina.
Si requiere una cotización, el personal la prepara. Una solicitud genera una
reserva pendiente; el sistema verifica los cupos dentro de una transacción y
registra los pasajeros. Cuando corresponde, se registra un anticipo con la
pasarela simulada y el personal lo valida antes de confirmar. Si vence el
plazo o se aprueba una cancelación, los cupos retenidos se liberan por
cálculo y el sistema comunica los cambios dentro de la propia plataforma.

Las reglas principales del proceso son cinco: no permitir sobreventa;
mantener una sola fuente de disponibilidad, calculada y no almacenada;
retener los cupos mediante una operación transaccional que bloquea la fila de
la salida; no confirmar automáticamente por el solo registro del pago; y
conservar la trazabilidad de las modificaciones, las cancelaciones y los
vencimientos.

\subsection{Diagrama de actividades con calles}

El diagrama con calles muestra las responsabilidades del cliente, del
personal, del administrador y del sistema durante la consulta, la solicitud,
la verificación y la confirmación de las reservas
(figura~\ref{fig:calles-negocio}).

\begin{figuraapa}{fig:calles-negocio}{Proceso de negocio por actor}
  \incluirDiagrama{procesos/mapa-procesos}
  \nota{Fuente: elaboración propia.}
\end{figuraapa}

\section{Desarrollo de la propuesta: metodología Kanban}

\subsection{Roles Kanban}

Las responsabilidades se distribuyeron entre los dos universitarios y un
representante de Vizcacha Tours, conforme a la tabla~\ref{tab:roles}.

\begin{tablaapa}{tab:roles}{Roles Kanban}
\small
\begin{tabular}{@{}p{3.4cm}p{4.6cm}p{6.6cm}@{}}
\toprule
Rol & Responsable & Funciones \\
\midrule
Propietario del servicio & Representante de Vizcacha Tours &
Priorizar necesidades, validar reglas y aceptar entregables. \\
Gestor de operaciones & Carmen Nelinda Ozuna Villca &
Coordinar el flujo, los límites de trabajo en curso, los bloqueos y las
revisiones. \\
Equipo de trabajo & Roberto Carlos Ereño Acuña y Erick Fernando Velasco
Flores & Analizar, diseñar, programar, probar, documentar y corregir. \\
\bottomrule
\end{tabular}
\nota{Elaboración propia.}
\end{tablaapa}

\subsection{Identificación de actores y responsabilidades}

Los actores del sistema y sus responsabilidades se resumen en la
tabla~\ref{tab:actores}. Representan roles del sistema, no necesariamente
personas distintas: la confirmación final, la validación de los pagos y la
ejecución de las cancelaciones corresponden al personal autorizado.

\begin{tablaapa}{tab:actores}{Actores y responsabilidades}
\small
\begin{tabular}{@{}p{3.6cm}p{11cm}@{}}
\toprule
Actor & Responsabilidad en el sistema \\
\midrule
Administrador & Gestiona paquetes, proveedores, salidas, fechas y
capacidades; administra los detalles de servicios de los proveedores. \\
Personal de agencia & Consulta la disponibilidad, registra clientes,
cotiza, administra reservas, valida pagos, confirma solicitudes y consulta
reportes. \\
Cliente / turista & Crea su cuenta, consulta paquetes, salidas y fichas de
salida, solicita reservas, registra pasajeros y pagos, solicita la
cancelación y consulta avisos. \\
\bottomrule
\end{tabular}
\nota{Elaboración propia.}
\end{tablaapa}

\subsection{Políticas de trabajo}

Cada tarjeta lleva identificador, responsable, criterios de aceptación y
fechas de registro. Las columnas del tablero tienen límites de trabajo en
curso (4 en «En curso», 2 en «En pruebas / QA» y 1 en «Listo para
producción») y aplican una política de extracción: el equipo toma tareas
nuevas solo cuando existe capacidad disponible. Las tareas terminadas pasan a
revisión y se publican cuando cumplen los criterios de su historia; los
bloqueos se marcan en el tablero y se revisan diariamente, y la aceptación
corresponde al propietario del servicio.

\subsection{Historias de usuario}

Se conserva el listado del equipo, con su orden, prioridad, riesgo,
estimación y número de tareas. El backlog quedó formado por 19 historias,
55 puntos de historia y 54 tareas (tablas~\ref{tab:hus-1}
y~\ref{tab:hus-2}).

\begin{tablaapa}{tab:hus-1}{Historias de usuario}
\small
\begin{tabular}{@{}p{1.15cm}p{5.2cm}p{3.5cm}p{2.1cm}p{2.1cm}@{}}
\toprule
ID & Historia & Actor / canal & Prioridad / riesgo & Pts / tareas \\
\midrule
HU-08 & Iniciar sesión y acceder según el rol & Personal, Cliente / Web &
Alta / Alto & 3 pts / 3 tareas \\
HU-15 & Abrir una cuenta de cliente y acceder a su información & Cliente /
Web & Media / Alto & 3 pts / 3 tareas \\
HU-01 & Gestión de paquetes turísticos & Administrador / Web & Alta / Medio
& 3 pts / 3 tareas \\
HU-02 & Registro de fechas de salida y cupos & Administrador / Web & Alta /
Alto & 2 pts / 2 tareas \\
HU-03 & Consulta de disponibilidad & Cliente, Personal / Web, Oficina &
Alta / Medio & 2 pts / 2 tareas \\
HU-04 & Registro y actualización de la ficha del cliente & Personal / Web,
Oficina & Media / Medio & 2 pts / 2 tareas \\
HU-10 & Registrar proveedores y vincularlos a los paquetes & Administrador /
Web & Media / Bajo & 2 pts / 2 tareas \\
HU-05 & Elaboración de cotizaciones & Personal / Web, Oficina & Baja / Medio
& 3 pts / 3 tareas \\
HU-06 & Registro de reservas con validación de cupos & Personal, Cliente /
Web, Oficina & Alta / Alto & 5 pts / 4 tareas \\
HU-12 & Registro de pasajeros y vinculación a una reserva & Personal,
Cliente / Web & Media / Medio & 3 pts / 2 tareas \\
\bottomrule
\end{tabular}
\nota{Elaboración propia a partir del backlog del proyecto.}
\end{tablaapa}

\begin{tablaapa}{tab:hus-2}{Historias de usuario (continuación)}
\small
\begin{tabular}{@{}p{1.15cm}p{5.2cm}p{3.5cm}p{2.1cm}p{2.1cm}@{}}
\toprule
ID & Historia & Actor / canal & Prioridad / riesgo & Pts / tareas \\
\midrule
HU-07 & Seguimiento y actualización de reservas & Personal / Web & Media /
Medio & 3 pts / 4 tareas \\
HU-17 & Confirmación de solicitudes y control de su vencimiento & Personal /
Web & Alta / Alto & 3 pts / 4 tareas \\
HU-09 & Modificar fechas o capacidad de una salida & Administrador / Web &
Media / Alto & 3 pts / 2 tareas \\
HU-11 & Consultar reportes de ocupación y reservas por fecha & Personal /
Web & Baja / Medio & 2 pts / 2 tareas \\
HU-13 & Solicitud y ejecución de la cancelación de una reserva & Cliente,
Personal / Web, Oficina & Media / Alto & 3 pts / 2 tareas \\
HU-14 & Registro y validación de anticipos y saldos de una reserva &
Cliente, Personal / Web, Oficina & Alta / Medio & 3 pts / 3 tareas \\
HU-16 & Avisos de cambios de fecha y capacidad de una salida & Cliente / Web
& Baja / Medio & 2 pts / 2 tareas \\
HU-18 & Ficha de la salida & Cliente, Personal / Web & Alta / Medio &
5 pts / 5 tareas \\
HU-19 & Detalles de servicios de proveedores & Administrador / Web & Media /
Medio & 3 pts / 4 tareas \\
\midrule
\multicolumn{4}{@{}p{12.05cm}}{Total: 19 historias, 55 puntos de historia y
54 tareas.} & \\
\bottomrule
\end{tabular}
\nota{Elaboración propia a partir del backlog del proyecto.}
\end{tablaapa}

\subsubsection{Criterios de aceptación}

Cada historia tiene criterios de aceptación verificables, apoyados en las
reglas de negocio del sistema (tablas~\ref{tab:criterios-1}
y~\ref{tab:criterios-2}).

\begin{tablaapa}{tab:criterios-1}{Criterios de aceptación}
\small
\begin{tabular}{@{}p{1.3cm}p{13cm}@{}}
\toprule
Historia & Criterio verificable \\
\midrule
HU-08 & Credenciales válidas y permisos según rol; el acceso no autorizado
se rechaza y queda registrado. \\
HU-15 & Cuenta única por ficha de cliente; el cliente solo consulta sus
propios datos. \\
HU-01 & Paquetes con código generado, datos obligatorios y estados borrador,
activo e inactivo; solo el activo se ofrece. \\
HU-02 & Salida con fecha válida, capacidad positiva y precio; el cupo se
calcula y un paquete no tiene dos salidas el mismo día. \\
HU-03 & La web y la oficina muestran el mismo cupo disponible, calculado con
las reservas pendientes y confirmadas vigentes. \\
HU-04 & Ficha con documento y correo únicos, sin duplicados; una ficha puede
existir sin cuenta de usuario. \\
HU-10 & Proveedores con tipos declarados; el vínculo con el paquete no se
repite y conserva el costo. \\
HU-05 & Cotización con número único, importe calculado según salida y
pasajeros, y plazo de validez; no descuenta cupo ni crea reserva. \\
HU-06 & Solicitud pendiente solo si hay cupos, dentro de una transacción con
bloqueo de fila; sin sobreventa y con vencimiento a 24 horas. \\
HU-12 & Pasajeros vinculados a la reserva; la cantidad registrada coincide
con la declarada y el documento no se repite. \\
\bottomrule
\end{tabular}
\nota{Elaboración propia.}
\end{tablaapa}

\begin{tablaapa}{tab:criterios-2}{Criterios de aceptación (continuación)}
\small
\begin{tabular}{@{}p{1.3cm}p{13cm}@{}}
\toprule
Historia & Criterio verificable \\
\midrule
HU-07 & Solo los estados permitidos cambian, con registro de cambios;
modificar la salida o la cantidad vuelve a verificar los cupos. \\
HU-17 & Confirmación solo del personal, con anticipo validado y manifiesto
completo; las solicitudes vencidas se rechazan y liberan el cupo. \\
HU-09 & La capacidad nueva no es menor a los pasajeros comprometidos; el
cambio queda auditado y genera aviso. \\
HU-11 & Ocupación y reservas filtradas por fecha o periodo, con las mismas
cifras que la consulta de disponibilidad. \\
HU-13 & El cliente solicita la cancelación; el personal la ejecuta con
motivo y fecha, y libera el cupo una sola vez. \\
HU-14 & Anticipo no menor al 30\% del monto; el pago web queda registrado
hasta que el personal lo valida y ningún pago supera el saldo. \\
HU-16 & Aviso interno al cliente afectado, con el valor anterior, el nuevo y
la fecha del cambio. \\
HU-18 & Ficha de la salida con mapa, itinerario, servicios y proveedores; el
cliente no ve costos internos. \\
HU-19 & Detalles por servicio del proveedor con título y orden; las
asignaciones operativas de la salida no tienen costos. \\
\bottomrule
\end{tabular}
\nota{Elaboración propia.}
\end{tablaapa}

\subsection{Diagrama de casos de uso}

Las 19 historias del backlog se agrupan en 15 casos de uso del sistema, de
modo que cada caso representa lo que el actor quiere conseguir y no las
tareas internas (figura~\ref{fig:casos-uso}). La consulta, la solicitud de
reserva y el pago pueden provenir del cliente con cuenta o de la atención del
personal, y en ambos casos terminan en el mismo proceso validado. El turista
solicita reservas y registra pagos; el personal conserva la confirmación, la
validación de los pagos y las cancelaciones; el administrador mantiene el
catálogo. Gestionar la reserva es el caso general que extienden la
confirmación, la cancelación y el registro de pagos: el registro del pago del
cliente no equivale a su validación por el personal. La ficha de la salida
(HU-18), con el mapa del recorrido, el itinerario y los detalles de los
servicios de los proveedores (HU-19), forma parte de la consulta de
disponibilidad.

\begin{landscape}
\thispagestyle{paisaje}
\begin{figuraapa}{fig:casos-uso}{Casos de uso del sistema}
  \incluirDiagrama{casos-uso/casos-uso}
  \nota{Fuente: elaboración propia.}
\end{figuraapa}
\end{landscape}

\subsection{Lista de backlog y tarjetas Kanban}

El flujo de trabajo comienza con la identificación de necesidades y su
registro en el backlog. Se seleccionan las historias considerando la
prioridad, el riesgo y las dependencias; cada tarjeta define sus criterios y
sus tareas antes de pasar a progreso, donde se realiza el análisis, el
diseño y la programación. En pruebas se verifica la funcionalidad y, una vez
comprobada y publicada, la tarjeta llega a hecho. Si falla, retorna a
progreso cuando exista capacidad.

\subsubsection{Backlog y orden de implementación}

El backlog se organizó en cinco bloques de implementación, con el orden
resumido en la tabla~\ref{tab:backlog}. La primera ventana de trabajo
construyó como producto mínimo viable el acceso por rol, el catálogo de
paquetes y la ficha del cliente (HU-08, HU-01 y HU-04), y las historias
HU-18 y HU-19, incorporadas después de las pruebas del ciclo, quedaron en
«Por hacer» al cierre del periodo.

\begin{tablaapa}{tab:backlog}{Backlog y orden de implementación}
\small
\begin{tabular}{@{}p{3.4cm}p{5.6cm}p{5.6cm}@{}}
\toprule
Bloque & Historias & Resultado esperado \\
\midrule
1. Acceso y catálogo & HU-08, HU-15, HU-01, HU-02, HU-03, HU-18 & Acceso,
paquetes, consulta de cupos y ficha de la salida. \\
2. Atención & HU-04, HU-10, HU-19, HU-05 & Clientes, proveedores, detalles
de servicios y cotizaciones. \\
3. Reserva & HU-06, HU-12, HU-07, HU-17 & Solicitud, pasajeros, seguimiento
y confirmación. \\
4. Cambios y pagos & HU-09, HU-13, HU-14 & Cambios, cancelación y validación
de pagos. \\
5. Información & HU-11, HU-16 & Reportes y avisos internos. \\
\bottomrule
\end{tabular}
\nota{Elaboración propia.}
\end{tablaapa}

Aunque HU-14 figura en el cuarto bloque, la confirmación de una reserva
requiere el pago validado, de modo que ninguna historia quedó por terminada
hasta integrar esa validación. El orden se ajustó en las reuniones de
reposición cuando las dependencias entre historias lo exigieron.

\subsubsection{Tablero Kanban y límites WIP}

El tablero de trabajo tiene cinco columnas con límites de trabajo en curso,
que son los que imponen la política de extracción y el ritmo del equipo
(figura~\ref{fig:flujo-kanban}). El estado del tablero al cierre de la
primera ventana, el 25 de noviembre de 2026, se resume en la
tabla~\ref{tab:tablero}: el backlog original quedó completo en «Hecho» y solo
las historias nuevas HU-18 y HU-19 permanecen en «Por hacer».

\begin{tablaapa}{tab:tablero}{Tablero Kanban al 25 de noviembre de 2026}
\small
\begin{tabular}{@{}p{4.6cm}p{1.6cm}p{7.8cm}@{}}
\toprule
Columna & WIP & Tarjetas al 25/11/2026 \\
\midrule
Por hacer & --- & 9: HU-18 (5 tareas) y HU-19 (4 tareas). \\
En curso & 4 & Sin tarjetas. \\
En pruebas / QA & 2 & Sin tarjetas. \\
Listo para producción & 1 & Sin tarjetas. \\
Hecho & --- & 45: las 17 historias del backlog original, incluidas las
tareas de ajuste detectadas en las pruebas. \\
\bottomrule
\end{tabular}
\nota{Elaboración propia a partir del tablero del proyecto.}
\end{tablaapa}

\begin{figuraapa}{fig:flujo-kanban}{Flujo de tarjetas Kanban}
  \incluirDiagrama{procesos/flujo-tablero-kanban}
  \nota{Fuente: elaboración propia.}
\end{figuraapa}

Con el ritmo observado, de 1,75 tareas por día de equipo, la ley de Little
relaciona el límite de trabajo en curso con el tiempo de ciclo: con cuatro
tarjetas en curso, cada tarea tarda en promedio 2,29 días en atravesar el
sistema. El equipo revisa semanalmente el tablero, mide el tiempo de ciclo y
la cantidad de historias terminadas, y ajusta los límites según la capacidad
real de los dos integrantes.

\subsubsection{Aplicación del método STATIK}

STATIK (Systems Thinking Approach to Introducing Kanban) es el procedimiento
de seis pasos para diseñar un sistema Kanban ajustado a un servicio
\parencite{kanbanGuia}. Se aplicó de manera iterativa, considerando la
demanda de Vizcacha Tours y la capacidad del equipo, conforme a la
tabla~\ref{tab:statik}.

\begin{tablaapa}{tab:statik}{Aplicación de STATIK}
\small
\begin{tabular}{@{}p{4.2cm}p{10cm}@{}}
\toprule
Paso & Aplicación al proyecto \\
\midrule
1. Identificar fuentes de insatisfacción & Diferencias de cupos entre
canales, demoras de confirmación y dificultades de coordinación
identificadas en el diagnóstico. \\
2. Analizar la demanda & 19 historias clasificadas por actor, canal,
prioridad, riesgo y dependencia. \\
3. Analizar la capacidad & Dos integrantes con 4 horas por día; ritmo medido
de 4,00 horas-persona por punto y 1,75 tareas por día. \\
4. Modelar el flujo & Cinco columnas: por hacer, en curso, en pruebas,
listo para producción y hecho. \\
5. Identificar clases de servicio & Estándar para las historias; urgente
para fallas críticas de cupos o seguridad, previa autorización. \\
6. Diseñar el sistema Kanban & Tablero con límites de 4, 2 y 1; política de
extracción; reglas de entrada y salida; reuniones de seguimiento. \\
\bottomrule
\end{tabular}
\nota{Elaboración propia con \textcite{kanbanGuia}.}
\end{tablaapa}

\subsubsection{Análisis} Las actividades de análisis y sus resultados se
resumen en la tabla~\ref{tab:analisis}.

\begin{tablaapa}{tab:analisis}{Actividades de análisis}
\small
\begin{tabular}{@{}p{7.2cm}p{7.2cm}@{}}
\toprule
Actividad & Resultado \\
\midrule
Entrevistas y revisión autorizada de registros & Necesidades y reglas del
proceso documentadas. \\
Precisar reglas de cupos, vencimientos y permisos & Catálogo de reglas de
negocio y criterios de aceptación. \\
Validar el alcance de HU-06, HU-14 y HU-17 & Estados y validaciones definidos
para los dos canales. \\
Revisar el boceto con la docente & Casos de uso y modelo de datos
ajustados. \\
\bottomrule
\end{tabular}
\nota{Elaboración propia.}
\end{tablaapa}

\subsubsection{Diseño} Las actividades de diseño y sus resultados se resumen
en la tabla~\ref{tab:diseno}.

\begin{tablaapa}{tab:diseno}{Actividades de diseño}
\small
\begin{tabular}{@{}p{5.2cm}p{4.6cm}p{4.6cm}@{}}
\toprule
Actividad & Desarrollo & Salida \\
\midrule
Arquitectura y prototipos & Modelar componentes, clases y pantallas & Diseño
aprobado \\
Datos y reglas & Definir claves, transacciones y estados & Esquema
relacional \\
Pruebas & Diseñar escenarios de concurrencia, permisos y vencimiento &
Casos de prueba \\
\bottomrule
\end{tabular}
\nota{Elaboración propia.}
\end{tablaapa}

\subsection{Diagrama de clases y objetos}

El modelo de clases representa las entidades del negocio y sus asociaciones
(figura~\ref{fig:clases}). El paquete pertenece al catálogo de lugares,
declara su itinerario día por día y sus servicios incluidos; la salida
pertenece a un paquete y mantiene la capacidad programada. Una reserva se
vincula a un cliente y a una salida; los pagos y los pasajeros se relacionan
con la reserva. Los proveedores declaran sus tipos de servicio y sus
detalles, se vinculan a los paquetes con su costo y se asignan a las salidas
como asignaciones operativas. La disponibilidad se calcula a partir de los
compromisos vigentes y no se almacena, al igual que el saldo pendiente, que
se calcula sobre los pagos validados.

\begin{landscape}
\thispagestyle{paisaje}
\begin{figuraapa}{fig:clases}{Diagrama de clases del sistema}
  \incluirDiagrama{estructural/diagrama-clases}
  \nota{Fuente: elaboración propia.}
\end{figuraapa}
\end{landscape}

\subsubsection{Diagrama de objetos}

El ejemplo muestra, en un instante concreto, una reserva pendiente de dos
pasajeros para una salida con capacidad de doce plazas, originada en una
cotización aceptada y con un anticipo registrado con la pasarela simulada que
aún no ha sido validado (figura~\ref{fig:objetos}). Los datos son ficticios
y no corresponden a registros reales de la agencia.

\begin{landscape}
\thispagestyle{paisaje}
\begin{figuraapa}{fig:objetos}{Diagrama de objetos: solicitud de reserva y anticipo}
  \incluirDiagrama{estructural/diagrama-objetos}
  \nota{Fuente: elaboración propia.}
\end{figuraapa}
\end{landscape}

\subsection{Diagramas del sistema}

\subsubsection{Diagrama de secuencia}

Las figuras~\ref{fig:secuencia-hu06} y~\ref{fig:secuencia-hu17} muestran dos
flujos representativos del sistema: el registro de una reserva con la
validación de cupos dentro de una transacción, por cualquiera de los dos
canales (HU-06), y la confirmación de las solicitudes con su control de
vencimiento (HU-17). El resto de los flujos del sistema sigue el mismo patrón
de interfaz, control, entidad y tabla, con el bloqueo de la fila de la
salida y el registro auditado de cada cambio.

\begin{figuraapa}{fig:secuencia-hu06}{Secuencia: registro de la reserva con validación de cupos (HU-06)}
  \incluirDiagrama{secuencia/hu-06-registrar-reserva}
  \nota{Fuente: elaboración propia.}
\end{figuraapa}

\begin{figuraapa}{fig:secuencia-hu17}{Secuencia: confirmación de solicitudes y vencimiento (HU-17)}
  \incluirDiagrama{secuencia/hu-17-confirmar-solicitudes}
  \nota{Fuente: elaboración propia.}
\end{figuraapa}

\subsubsection{Diagrama de actividades}

La figura~\ref{fig:estados-reserva} resume el ciclo de vida de una reserva y
las condiciones que gobiernan cada transición: el plazo de 24 horas, el
anticipo validado, el manifiesto completo y la ejecución del personal.

\begin{figuraapa}{fig:estados-reserva}{Gestión de los estados de la reserva}
  \incluirDiagrama{procesos/estados-reserva}
  \nota{Fuente: elaboración propia.}
\end{figuraapa}

\subsubsection{Diagrama de componentes}

El sistema se organiza en tres capas (figura~\ref{fig:componentes}): la
interfaz web, que llama a la API del servidor y a la autenticación; el
servidor, con sus rutas, validaciones y servicios; y la base de datos, donde
las reglas críticas de cupos, pagos y confirmación viven en funciones
almacenadas que la aplicación invoca sin reimplementarlas.

\begin{figuraapa}{fig:componentes}{Componentes del sistema}
  \incluirDiagrama{estructural/diagrama-componentes}
  \nota{Fuente: elaboración propia.}
\end{figuraapa}

\subsubsection{Diagrama de despliegue}

El despliegue propuesto (figura~\ref{fig:despliegue}) distribuye la interfaz
y las funciones del servidor en una plataforma en la nube, con la base de
datos PostgreSQL administrada en Neon y accesos diferenciados: la aplicación
opera mediante una conexión con agrupador y el equipo de desarrollo aplica el
esquema por conexión directa. El cobro se emula con la pasarela simulada y
los avisos se muestran dentro de las pantallas del sistema.

\begin{figuraapa}{fig:despliegue}{Despliegue del sistema}
  \incluirDiagrama{estructural/diagrama-despliegue}
  \nota{Fuente: elaboración propia.}
\end{figuraapa}

\subsection{Validación del sistema}

Las reglas de negocio se verificaron con pruebas automatizadas sobre la base
de datos, que cubren la gestión de paquetes y salidas, la validación de
cupos, los anticipos, la confirmación, el vencimiento, la modificación, la
cancelación y la finalización de reservas; cada prueba termina en
\texttt{ROLLBACK}, de modo que no altera los datos. Además se ejecutaron
pruebas de API y de navegador para el acceso por rol y para los flujos de
consulta, y las tareas atravesaron la columna de pruebas del tablero antes de
publicarse. Entre los escenarios comprobados está la reserva concurrente del
último cupo, que la transacción con bloqueo de fila resuelve sin sobreventa.

Quedan pendientes de validación las historias HU-18 y HU-19, incorporadas al
alcance después de las pruebas del ciclo. Los indicadores previstos para la
próxima revisión son el tiempo de confirmación, las discrepancias de cupos y
la cantidad de solicitudes atendidas.
